\documentclass[11pt]{article}
\ifdefined\pdfinfoomitdate
  \pdfinfoomitdate=1
\fi
\ifdefined\pdftrailerid
  \pdftrailerid{}
\fi
\ifdefined\pdfsuppressptexinfo
  \pdfsuppressptexinfo=15
\fi

\usepackage[a4paper,margin=1in]{geometry}
\usepackage{amsmath}
\usepackage{amssymb}
\usepackage{siunitx}

\title{Cosmic-Ray Threshold-Rate Reference and Approximation}
\author{}
\date{}

\begin{document}
\maketitle

\section*{Scope and data convention}

The bundled \texttt{data.csv} is interpreted as
\[
E \; [\mathrm{eV}], \qquad
F(E) \; [(\mathrm{m}^2\,\mathrm{sr}\,\mathrm{s}\,\mathrm{GeV})^{-1}] .
\]
The values were digitized from a plotted spectrum rather than supplied as a
primary experimental dataset. No point-by-point uncertainties or reproducible
original digitization log were included. Consequently, the numerical results
below are reference calculations for this CSV, not exact physical predictions,
and the displayed precision is intentionally limited. The source graphic is
attributed to Sven Lafebre (8 January 2019); this project uses the
Creative Commons Attribution-ShareAlike 3.0 option offered on the source file
page for the local raster and copyrightable adaptations. See
\texttt{DATA\_PROVENANCE.md} and \texttt{NOTICE} for full details.

\section*{Explicitly truncated CSV reference}

Let $E_{\max}$ be the largest tabulated energy, approximately
$3.36\times10^{20}\,\mathrm{eV}$. For a one-sided horizontal surface of area
$A$ exposed to an isotropic downward intensity, the ideal crossing-rate
reference used by the code is
\[
R_{\mathrm{CSV}}(>E_{\mathrm{th}};E_{\max})
=
\pi A_{\mathrm{m}^2}\,10^{-9}
\int_{E_{\mathrm{th}}}^{E_{\max}} F(E)\,dE_{\mathrm{eV}} .
\]
For $A=\SI{1}{km^2}=10^6\,\si{m^2}$ this becomes
\[
R_{\mathrm{CSV}}(>E_{\mathrm{th}};E_{\max})
=
\pi\times10^{-3}
\int_{E_{\mathrm{th}}}^{E_{\max}} F(E)\,dE_{\mathrm{eV}} .
\]
The upper limit is finite by definition. The calculation does \emph{not}
assert that the physical flux is zero above $E_{\max}$.

The code also offers an explicit \texttt{last-segment} option. If the final
log-log segment is extrapolated as
$F(E)=F(E_{\max})(E/E_{\max})^{\gamma_{\mathrm{last}}}$ with
$\gamma_{\mathrm{last}}<-1$, its model-dependent tail contribution is
\[
\int_{E_{\max}}^{\infty} F(E)\,dE
=
\frac{F(E_{\max})E_{\max}}{-\left(\gamma_{\mathrm{last}}+1\right)} .
\]
This extrapolation can materially change results near the top of the CSV and is
therefore never applied silently.

\section*{Acceptance assumptions}

The factor $\pi A$ is the projected geometric acceptance of an ideal, one-sided,
horizontal surface for a downward isotropic intensity. It is not a complete
instrument response. Detector efficiency, dead time, trigger and reconstruction
response, atmospheric shower development, and energy- or direction-dependent
effective area are not included. The Python interface can instead accept a
caller-supplied scalar aperture in $\si{m^2.sr}$, but realistic instruments
usually require integrating a full energy- and direction-dependent response.

\section*{Piecewise approximation}

The following approximation is calibrated to the bundled CSV, the finite
$E_{\max}$ truncation convention, a flat horizontal $\SI{1}{km^2}$ acceptance,
and the range
\[
10^{10}\,\mathrm{eV}\le E\le10^{20}\,\mathrm{eV}.
\]
Its units are $\mathrm{s}^{-1}\,\mathrm{km}^{-2}$. The coefficients and
validation statistics are generated by the companion Python script from the
same implementation used by the command-line program.

% Generated by generate_artifacts.py; do not edit by hand.
\[
R_{\mathrm{fit}}(>E) = R_i\left(\frac{E}{E_i}\right)^{-p_i},
\qquad E_i\le E<E_{i+1}.
\]
\begin{center}
\begin{tabular}{c c c c}
\hline
$E_i$ [eV] & $E_{i+1}$ [eV] & $R_i$ [$\mathrm{s}^{-1}\,\mathrm{km}^{-2}$] & $p_i$ \\
\hline
$10^{10}$ & $4.5 \times 10^{11}$ & $3.898 \times 10^{8}$ & $1.516$ \\
$4.5 \times 10^{11}$ & $6.4 \times 10^{15}$ & $1.216 \times 10^{6}$ & $1.672$ \\
$6.4 \times 10^{15}$ & $6 \times 10^{18}$ & $1.382 \times 10^{-1}$ & $2.148$ \\
$6 \times 10^{18}$ & $3.3 \times 10^{19}$ & $5.703 \times 10^{-8}$ & $1.880$ \\
$3.3 \times 10^{19}$ & $10^{20}$ & $2.313 \times 10^{-9}$ & $2.750$ \\
\hline
\end{tabular}
\end{center}
The final interval includes its upper endpoint. Table coefficients are rounded.
For the unrounded fit, 10,001 thresholds uniformly spaced in $\log_{10} E$ give sampled absolute relative errors of $5.33\%$ (median), $14.58\%$ (95th percentile), and $17.03\%$ (maximum), relative to the CSV reference.


\section*{Selected thresholds}

The table uses integer-decade thresholds and modest output precision. The
``CSV-truncated reference'' column integrates only to the final tabulated point.
Values near $10^{20}\,\mathrm{eV}$ should therefore be treated as sensitive to
the unspecified high-energy tail.

% Generated by generate_artifacts.py; do not edit by hand.
\begin{center}
\small
\begin{tabular}{c c c}
\hline
Threshold $E_{\mathrm{th}}$ [eV] & Piecewise fit [$\mathrm{s}^{-1}\,\mathrm{km}^{-2}$] & CSV-truncated reference [$\mathrm{s}^{-1}\,\mathrm{km}^{-2}$] \\
\hline
$10^{12}$ & $3.298 \times 10^{5}$ & $3.298 \times 10^{5}$ \\
$10^{13}$ & $7.085 \times 10^{3}$ & $6.031 \times 10^{3}$ \\
$10^{14}$ & $1.522 \times 10^{2}$ & $1.519 \times 10^{2}$ \\
$10^{15}$ & $3.271$ & $3.474$ \\
$10^{16}$ & $4.387 \times 10^{-2}$ & $4.953 \times 10^{-2}$ \\
$10^{17}$ & $4.04 \times 10^{-4}$ & $4.332 \times 10^{-4}$ \\
$10^{18}$ & $2.49 \times 10^{-6}$ & $2.406 \times 10^{-6}$ \\
$10^{19}$ & $2.177 \times 10^{-8}$ & $2.225 \times 10^{-8}$ \\
$10^{20}$ & $1.044 \times 10^{-10}$ & $1.044 \times 10^{-10}$ \\
\hline
\end{tabular}
\end{center}


For an ideal flat surface of another area $A$ in $\si{km^2}$, both columns scale
linearly with $A$. The documented fit is not automatically applied to modified
CSV data or another tail convention; it must be explicitly recalibrated.

\end{document}
